\documentclass[10pt,a4paper]{article}
\usepackage{fontspec,amsmath,amssymb,graphicx}
\usepackage[margin=18mm]{geometry}
\setmainfont{Times New Roman}
\setlength{\parindent}{0pt}
\setlength{\parskip}{5pt}
\allowdisplaybreaks
\sloppy
\emergencystretch=2em
\begin{document}
\section*{BT\_3\_6\_3\_06 — Lời giải bổ sung}
\subsection*{Kiến thức và trình tự giải}
\textbf{Bước 1 — Chọn Kirchhoff (3D) hoặc Poisson (2D).} Với li độ đầu bằng không, vận tốc $\psi$, tốc độ $c$, đặt $\rho=ct$. Trong ba chiều,


\[
u(P,t)=\frac t{4\pi}\int_{S^2}\psi(P+\rho\omega)d\omega.
\]


Trong hai chiều,


\[
u(P,t)=\frac1{2\pi c}\int_{|Q-P|<ct}\frac{\psi(Q)}{\sqrt{c^2t^2-|Q-P|^2}}dQ.
\]


Với li độ $f$ khác không, thêm đạo hàm theo $t$ của chính biểu thức tương ứng khi thay $\psi$ bằng $f$. Với nguồn và dữ liệu đầu không, tích phân theo thời gian của nghiệm có vận tốc là nguồn tại từng thời điểm (Duhamel).

\textbf{Bước 2 — Tận dụng hình học và đối xứng.} Chỉ thị miền $D$ biến Kirchhoff thành $t$ nhân tỷ lệ diện tích mặt cầu nằm trong $D$. Chọn trục qua tâm miền và $P$, đặt $\mu=\cos\theta$, $d\omega=d\mu\,d\phi$. Nếu điều kiện không phụ thuộc $\phi$, tỷ lệ là $|E|/2$ với $E\subset[-1,1]$. Mỗi nửa không gian cho một bất đẳng thức tuyến tính theo $\mu$; quả cầu cho một ngưỡng chỏm. Miền hợp/rút bỏ cho tổng/hiệu nếu các phần không chồng diện tích.

\textbf{Bước 3 — Viết chính xác các khoảng hợp lệ.} Giải bất đẳng thức gốc, sau đó mới bình phương khi có căn. Chia $[-1,1]$ ở nghiệm phương trình biên, kiểm tra dấu ở trung điểm từng khoảng, loại nghiệm ngoại lai. Nếu còn góc phương vị, dùng độ dài tập $\cos\phi<v$ rồi còn một tích phân theo $\mu$. Với Poisson, tích phân theo một tọa độ cho hiệu arcsin; cắt các cận vào đĩa trước khi tính.

\textbf{Bước 4 — Tính nguồn bằng nguyên hàm thời gian.} Với quả cầu bán kính $R$ cách $P$ khoảng $r>0$, đặt


\[
K_R(r,t)=t\left[\frac{R^2-(r-ct)^2}{4rct}\right]_0^1,\quad J_R(r,t)=\int_0^tK_R(r,s)ds.
\]


Tách ba giai đoạn: cả mặt cầu nằm trong miền, giao một chỏm, không giao. Nguồn bật trên $[0,T]$ cho $J_R(r,t)-J_R(r,(t-T)_+)$. Trường hợp $r=0$ tính trực tiếp, tránh chia cho $r$.

\textbf{Bước 5 — Kiểm tra chuẩn hóa, mặt sóng và cực đại.} Nếu $\psi=1$ trên toàn không gian thì $u=t$, dùng để kiểm tra hệ số. Với dữ liệu bước nhảy, chỉ kết luận cổ điển ngoài các mặt sóng. Cực đại phải xét cả điểm chuyển nhánh và giá trị một phía; nêu rõ supremum nếu chỉ đạt theo giới hạn. Nếu dùng tích phân số, ghi tham số và cách kiểm tra hội tụ; đồ thị không tự chứng minh một cực đại toàn cục.
\subsection*{Lời giải áp dụng cho bài}
\textbf{Bước 1. Chồng nghiệm vận tốc và nguồn.} Hai cầu vận tốc có dấu trái nhau cho $K_1(d_+,t)-K_1(d_-,t)$. Nguồn chỉ bật khi $0<s<T$, nên phần nguồn là $\int_0^{\min(t,T)}K_1(r,t-s)ds$. Đổi $\tau=t-s$ cho hiệu hai nguyên hàm $J$.

\textbf{Bước 2. Dùng đối xứng và thời gian yêu cầu.} Ở $x=0$, khoảng cách tới hai tâm $(\pm2,0,0)$ bằng nhau, nên phần vận tốc triệt tiêu hoàn toàn; câu (a) chỉ còn nguồn với $T=1$. Câu (b) giữ hai khoảng cách khác nhau và thay $t=r$ vào biểu thức, không thay $d_\pm$ bằng $r$.

\textbf{Bước 3. Công thức áp dụng.}

Gọi $d_\pm=|P-(\pm2,0,0)|$, $r=|P|$. Chồng Kirchhoff và Duhamel cho
\[u=K_1(d_+,t)-K_1(d_-,t)+J_1(r,t)-J_1(r,(t-T)_+),\]
với $c=1$, $K,J$ có biểu thức chỏm và nguyên hàm ở hai bài trước.
(a) $P=(0,y,z)$ cho $d_+=d_-$, nên $u=J_1(\sqrt{y^2+z^2},t)-J_1(\sqrt{y^2+z^2},(t-1)_+)$.
(b) Thay $t=r$: $u=K_1(d_+,r)-K_1(d_-,r)+J_1(r,r)-J_1(r,(r-T)_+)$. Đây là biểu thức tường minh bằng min/max và đa thức, áp dụng mọi $P$ kể cả $r=0$ theo giới hạn.

\textbf{Bước 4. Đối chiếu và kết luận.} Tại $t=0$, cả $J$ và $K$ bằng $0$; đạo hàm của phần $K$ cho đúng hiệu hai chỉ thị cầu và phần $J$ có vận tốc đầu không. Với $t=r=0$, kết quả là $0$ và dùng nhánh tâm của $J$ để tránh chia $r$. Các min/max của $K,J$ nêu trong phần kiến thức và Ví dụ 3.6.2.5 làm biểu thức câu (b) hoàn toàn tường minh.

\textbf{Chi tiết các hàm dùng trong biểu thức tường minh.} Với tốc độ $1$, bán kính $1$, $d>0,t>0$,
\[K_1(d,t)=t\left[\frac{1-(d-t)^2}{4dt}\right]_0^1.\]
Với $r>0$, đặt $b=|1-r|$, $q=[t]_b^{r+1}$ và $A(s)=((1-r^2)s+rs^2-s^3/3)/(4r)$:
\[J_1(r,t)=\tfrac12\min(t,(1-r)_+)^2+A(q)-A(b).\]
Tại tâm, $K_1(0,t)=t$ khi $t<1$, bằng $0$ khi $t>1$, $J_1(0,t)=\min(t,1)^2/2$; tại $t=0$, $K_1=J_1=0$. Khi đổi $\tau=t-s$, $s$ chạy $0$ đến $\min(t,T)$ nên $\tau$ chạy ngược từ $t$ đến $(t-T)_+$:
\[\int_0^{\min(t,T)}K_1(r,t-s)ds=\int_{(t-T)_+}^{t}K_1(r,\tau)d\tau=J_1(r,t)-J_1(r,(t-T)_+).\]\end{document}